% LaTeX companion of hw03.typ. Both formats are editable.
\section{Homework 3 --- submitted work}

\begin{qlnote}{Note}

Visual transcription of the personal finished submission. Source: \texttt{217-Hw-3-finished.pdf}; page references below are PDF pages.

\end{qlnote}

\subsection{Source p. 1 --- Exercise 20}

Reflection about the \(x - z\) plane sends \(\begin{pmatrix}
v_{1} \\
v_{2} \\
v_{3}
\end{pmatrix}\) to \(\begin{pmatrix}
v_{1} \\
{- v_{2}} \\
v_{3}
\end{pmatrix}\). The submitted standard matrix is \(\begin{pmatrix}
1 & 0 & 0 \\
0 & {- 1} & 0 \\
0 & 0 & 1
\end{pmatrix}\), accompanied by a sketch of the reflected vector.

\subsection{Source p. 2 --- Exercise 38(a)--(c)}

(a) Projection onto an arbitrary unit vector \(u = \begin{pmatrix}
u_{1} \\
u_{2}
\end{pmatrix}\) has matrix \(\begin{pmatrix}
u_{1}^{2} & {u_{1}u_{2}} \\
{u_{1}u_{2}} & u_{2}^{2}
\end{pmatrix}\). Its determinant is \(0\), so it is not invertible.

(b) A reflection matrix is \(\begin{pmatrix}
a & b \\
b & {- a}
\end{pmatrix}\) with \(a^{2} + b^{2} = 1\); its determinant is \(- 1\), so it is invertible.

(c) A rotation matrix is \(\begin{pmatrix}
a & {- b} \\
b & a
\end{pmatrix}\) with \(a^{2} + b^{2} = 1\); its determinant is \(a^{2} + b^{2} = 1\), so it is invertible.

\subsection{Source p. 3 --- Exercise 38(d); Exercise 18 begins}

(d) The horizontal and vertical shear matrices are \(\begin{pmatrix}
1 & k \\
0 & 1
\end{pmatrix}\) and \(\begin{pmatrix}
1 & 0 \\
k & 1
\end{pmatrix}\). Each has determinant \(1\), so each is invertible.

For Exercise 18, let \(A = \begin{pmatrix}
2 & 3 \\
{- 3} & 2
\end{pmatrix}\) and \(B = \begin{pmatrix}
a & b \\
c & d
\end{pmatrix}\). Equating \(AB\) and \(BA\) yields \(b = - c\) and \(a = d\), so \(B = \begin{pmatrix}
a & b \\
{- b} & a
\end{pmatrix}\).

\subsection{Source p. 4 --- Exercise 34}

For \(A = \begin{pmatrix}
1 & 1 \\
0 & 1
\end{pmatrix}\), the submitted computations are

\(A^{2} = \begin{pmatrix}
1 & 2 \\
0 & 1
\end{pmatrix}\), \(A^{3} = \begin{pmatrix}
1 & 3 \\
0 & 1
\end{pmatrix}\), \(A^{4} = \begin{pmatrix}
1 & 4 \\
0 & 1
\end{pmatrix}\), and \(A^{1001} = \begin{pmatrix}
1 & 1001 \\
0 & 1
\end{pmatrix}\).

Thus \(A^{k} = \begin{pmatrix}
1 & k \\
0 & 1
\end{pmatrix}\) for positive \(k\). Geometrically, repeated application is a horizontal shear by one unit each time.

\subsection{Source p. 5 --- Exercise 12}

For \(A = \begin{pmatrix}
2 & 5 & 0 & 0 \\
1 & 3 & 0 & 0 \\
0 & 0 & 1 & 2 \\
0 & 0 & 2 & 5
\end{pmatrix}\), the submission applies the invertible-matrix test through row reduction of \(\left\lbrack {A:I_{4}} \right\rbrack\). The displayed operations first exchange the first two rows and then use pivots to clear the two \(2 \times 2\) blocks.

\subsection{Source p. 6 --- Exercise 12; Exercise 34 begins}

The final augmented matrix is \(\left\lbrack {I_{4}:A^{-}1} \right\rbrack\), with

\(A^{-}1 = \begin{pmatrix}
3 & {- 5} & 0 & 0 \\
{- 1} & 2 & 0 & 0 \\
0 & 0 & 5 & {- 2} \\
0 & 0 & {- 2} & 1
\end{pmatrix}\).

By Theorem 2.4.5, \(A\) is invertible.

For a diagonal \(A = \begin{pmatrix}
a & 0 & 0 \\
0 & b & 0 \\
0 & 0 & c
\end{pmatrix}\), the submitted answer says \(A\) is invertible precisely when \(a,b,c \neq 0\), and then \(A^{-}1 = \begin{pmatrix}
\frac{1}{a} & 0 & 0 \\
0 & \frac{1}{b} & 0 \\
0 & 0 & \frac{1}{c}
\end{pmatrix}\). A diagonal matrix of arbitrary size is invertible iff every diagonal element is nonzero.

\subsection{Source p. 7 --- Part B, Problem 1(a)--(b)}

The source defines trace, determinant, transpose, and symmetric matrix, then asks truth values for transpose claims. The submission marks (a) false and takes \(A = \begin{pmatrix}
1 & 2 \\
2 & 1
\end{pmatrix}\) and \(B = \begin{pmatrix}
1 & 2 \\
1 & 2
\end{pmatrix}\). It computes \(AB = \begin{pmatrix}
3 & 6 \\
3 & 6
\end{pmatrix}\), \(B^{T} = \begin{pmatrix}
1 & 1 \\
2 & 2
\end{pmatrix}\), \(A^{T} = \begin{pmatrix}
1 & 2 \\
2 & 1
\end{pmatrix}\), and \(A^{T}B^{T} = \begin{pmatrix}
5 & 5 \\
4 & 4
\end{pmatrix}\), so \(\left( {AB} \right)^{T} \neq A^{T}B^{T}\).

For (b), it marks false using the same matrices and records \(AB = A^{T}B^{T} = \begin{pmatrix}
5 & 4 \\
4 & 5
\end{pmatrix}\) as a counterexample to the universal inequality.

\subsection{Source p. 8 --- Part B, Problem 1(c)}

The statement \(\left( {AB} \right)^{T} = B^{T}A^{T}\) is marked true. Let \(A\) be \(n \times p\) and \(B\) be \(p \times m\). The work writes both matrices by entries and notes that \(\left( {B^{T}A^{T}} \right)\) has the same \(m \times n\) shape as \(\left( {AB} \right)^{T}\).

\subsection{Source p. 9 --- Part B, Problem 1(c)--(d)}

The \(j,i\) entry of \(B^{T}A^{T}\) is computed as

\(\sum_{k = 1}^{p}b_{kj}a_{ik} = \sum_{k = 1}^{p}a_{ik}b_{kj}\),

the \(j,i\) entry of \(\left( {AB} \right)^{T}\). Since \(i,j\) are arbitrary, \(\left( {AB} \right)^{T} = B^{T}A^{T}\).

For (d), the submission begins induction: for a symmetric \(A = \begin{pmatrix}
a & b \\
b & a
\end{pmatrix}\), \(A^{1} = A\) is symmetric, and if \(A^{n} = \begin{pmatrix}
c & d \\
d & c
\end{pmatrix}\), then \(A^{n + 1} = A^{n}A\).

\subsection{Source p. 10 --- Part B, Problem 1(d)--(e)}

The inductive multiplication is

\(\begin{pmatrix}
c & d \\
d & c
\end{pmatrix}\begin{pmatrix}
a & b \\
b & a
\end{pmatrix} = \begin{pmatrix}
{ac + bd} & {ad + bc} \\
{ad + bc} & {ac + bd}
\end{pmatrix}\),

which is symmetric. Hence \(A^{n}\) is symmetric for all \(n \in \mathbb{N}\).

(e) is false: \(A = \begin{pmatrix}
1 & 3 \\
2 & 7
\end{pmatrix}\) is not symmetric, while \(A^{2} = \begin{pmatrix}
7 & 24 \\
16 & 55
\end{pmatrix}\) is recorded as symmetric in the submitted counterexample.

\subsection{Source p. 11 --- Part B, Problem 2(a)--(c)}

(a) True. A \(3 \times 3\) matrix with a zero row has at most two leading 1s, so \(\text{rank}(A) \leq 2\); by Theorem 2.4.3 it is not invertible.

(b) False. The counterexample is \(\begin{pmatrix}
1 & 1 & 0 \\
1 & 1 & 0 \\
0 & 0 & 1
\end{pmatrix}\), whose RREF has rank \(2 \neq 3\), so it is not invertible.

(c) True. If \(A\) is an invertible \(n \times n\) matrix, then \(A^{-}1A = \mathbb{A}^{-}1 = I_{n}\); therefore \(A^{-}1\) is invertible.

\subsection{Source p. 12 --- Part B, Problem 2(d)}

(d) True. If \(A\) is invertible, then \(\left( {A^{-}1} \right)^{n}\) is an inverse for \(A^{n}\). By associativity, the submission writes the products until adjacent \(A^{-}1A\) pairs become \(I_{n}\), so both products equal \(I_{n}\).

\subsection{Source p. 13 --- Part B, Problem 3(a)}

If there is an \(n \times m\) matrix \(B\) with \(BA = I_{n}\), suppose \(x_{1}\) and \(x_{2}\) solve \(Ax = 0\). Then \(A\left( {x_{1} - x_{2}} \right) = 0\). Multiplying on the left by \(B\) gives \(BA\left( {x_{1} - x_{2}} \right) = B0 = 0\); hence \(x_{1} - x_{2} = 0\) and \(x_{1} = x_{2}\). Therefore \(Ax = 0\) has a unique solution.

\subsection{Source p. 14 --- Part B, Problem 4(a)}

Statement (a) is marked true. Write \(A = \left\lbrack {v_{1},v_{2},\ldots,v_{m}} \right\rbrack\) and \(B = \left\lbrack {w_{1},w_{2},\ldots,w_{k}} \right\rbrack\). By Theorem 2.3.2,

\(AB = A\left\lbrack {w_{1},w_{2},\ldots,w_{k}} \right\rbrack = \left\lbrack {Aw_{1},Aw_{2},\ldots,Aw_{k}} \right\rbrack\).

If \(w_{i} = \begin{pmatrix}
w_{1i} \\
\ldots \\
w_{mi}
\end{pmatrix}\), then by Theorem 1.3.8,

\(Aw_{i} = w_{1i}v_{1} + w_{2i}v_{2} + \ldots + w_{mi}v_{m}\).

Thus every column of \(AB\) is a linear combination of the columns of \(A\).

\subsection{Source p. 15 --- Part B, Problem 4(b)}

Statement (b) is false. The counterexample uses \(A = \begin{pmatrix}
1 & 0 \\
0 & 0
\end{pmatrix}\) and \(B = \begin{pmatrix}
1 & 1 \\
1 & 1
\end{pmatrix}\), so \(AB = \begin{pmatrix}
1 & 1 \\
0 & 0
\end{pmatrix}\). Its first column \(\begin{pmatrix}
1 \\
0
\end{pmatrix}\) is not a linear combination of the columns \(\begin{pmatrix}
1 \\
1
\end{pmatrix}\) and \(\begin{pmatrix}
1 \\
1
\end{pmatrix}\) of \(B\): the submitted RREF shows that the corresponding system is inconsistent.

\subsection{Source p. 16 --- Part B, Problem 5(a)}

For a linear \(f:\mathbb{R}^{d}\rightarrow\mathbb{R}^{d}\), the proof by induction has base \(n = 1\), where \(f^{1}(x) = f\left( {f^{0}(x)} \right) = f(x)\) is linear. Assuming \(f^{n}\) is linear, write \(f(x) = Ax\). The key theorem gives a matrix \(B\) with \(f^{n{(x)}} = Bx\); hence

\(f^{n + 1}(x) = f\left( f^{n{(x)}} \right) = f\left( {Bx} \right) = A\left( {Bx} \right) = \left( {AB} \right)x\),

so it is linear. Thus \(f^{n}\) is linear for all \(n \in \mathbb{N}\).

\subsection{Source p. 17 --- Part B, Problem 5(a)}

The submission restates the inductive conclusion: the base case and inductive step prove that if \(f\) is a linear transformation, then \(f^{n}\) is a linear transformation for all \(n \in \mathbb{N}\).

\subsection{Source p. 18 --- Part B, Problem 5(b)}

Define \(f:\mathbb{R}^{2}\rightarrow\mathbb{R}^{2}\) by \(f(x) = - x\) if \(\left\| x \right\| = 2\), and \(f(x) = x\) otherwise. The submission says \(f\) is not linear: choose \(x_{1},x_{2}\) with \(\left\| x_{1} \right\| = 2\) and \(\left\| x_{2} \right\| \neq 2\), then \(f\left( {x_{1} + x_{2}} \right) = x_{1} + x_{2}\) while \(f\left( x_{1} \right) + f\left( x_{2} \right) = - x_{1} + x_{2}\). But \(f^{2}(x) = x\) for all \(x\), so \(f^{2}\) is the identity and is linear.

\subsection{Source p. 19 --- Part B, Problem 5(c)}

Let \(f(x) = Ax\). If \(f(x) = 0\) has a unique solution, then \(Ax = 0\) has a unique solution; by Theorem 1.3.4, \(\text{rank}(A) = d\), hence by Theorem 2.4.3 \(A\) is invertible. Problem 3 has shown that \(A^{n}\) is invertible for every \(n \in \mathbb{N}\). Since \(f^{n{(x)}} = A^{n}x\), Theorem 1.3.4 gives a unique solution to \(A^{n}x = 0\), proving the claim.
